\documentclass[11pt]{article}
\usepackage[utf8]{inputenc}
\usepackage{amsmath}
\usepackage{booktabs}

% simple isotope helper (avoids the isotope package dependency)
\newcommand{\iso}[2]{\textsuperscript{#1}#2}

\begin{document}

\begin{center}
  {\large\textbf{Offline PET/CT for proton range verification}}\\[2pt]
  {\large\textbf{summary of Parodi et al.\ (2007) and Knopf et al.\ (2009)}}
\end{center}
\medskip

\noindent
This note summarises the two foundational Massachusetts General Hospital (MGH)
studies of \emph{offline} PET/CT range verification of proton therapy: the
controlled phantom feasibility study of Parodi et al.~\cite{parodi2007a}
(Sec.~1), and the systematic patient-population study of the method's biological
and physical limitations by Knopf et al.~\cite{knopf2009} (Secs.~2--7). In both,
the measured $\beta^+$-activity distribution --- acquired on a commercial Biograph
PET/CT tens of minutes after irradiation --- is compared with a Monte Carlo
(FLUKA/Geant4) prediction built from the treatment plan, the CT and the
irradiation/imaging time course, and the proton range is read from the distal
fall-off of that activity. The phantom study establishes, under known geometry
and composition, that the method reaches millimetre range accuracy; the patient
study then shows that in vivo this accuracy is realised only in
well-co-registered bony structures of head-and-neck patients and fails at the
majority of positions and sites.

\subsection*{1. Phantom feasibility (Parodi et al.\ 2007)}

Parodi et al.~\cite{parodi2007a} irradiated plastic phantoms with SOBP proton
fields (15--16~cm water-equivalent range, 10~cm modulation, $8\times8$~cm$^2$
aperture) from the MGH 230~MeV cyclotron and imaged them 14--20~min later on an
LSO-based Biograph Sensation~16 in list mode, comparing with FLUKA/Geant4
predictions convolved with a 7~mm (transaxial) / 7.5~mm (axial) Gaussian. Four
configurations were used: a homogeneous PMMA cylinder (8~Gy single field, and
8~Gy two orthogonal fields), an inhomogeneous PMMA/bone/lung stack (2~Gy), and
PMMA with two titanium rods mimicking spinal implants (8~Gy). Three results
carry over to the present work.

\emph{Range accuracy.} The 50\% distal fall-off of the measured activity agreed
with the simulation \textbf{within 1~mm in every configuration} --- homogeneous
PMMA, the inhomogeneous stack, and even through the titanium implants --- and
held to within 1~mm when only the last 20~min was reconstructed (the
$\sim$2~Gy-equivalent low-statistics case, activity reduced by $2.29\times$).
Lateral field position was likewise verified to $<1$~mm. For the overlapping
two-field case the analysis had to be restricted to the 30--10\% fall-off to
exclude the second beam, and even there stayed within 1~mm.

\emph{Isotope content and yield.} A maximum \iso{11}{C} production of
$\sim$870~Bq\,ml$^{-1}$\,Gy$^{-1}$ was measured in PMMA after $\sim$6~min
irradiation. At the start of imaging the offline signal is $94$--$97\%$
\iso{11}{C}, $3$--$5\%$ \iso{13}{N}, with only trace \iso{15}{O} in the first
minutes --- the long-lived-isotope regime that offline acquisition selects.
Rescaled to tissue by carbon content, this corresponds to $\sim$0.3
(muscle)--1.4 (bone/adipose)~kBq\,ml$^{-1}$ after a 2~Gy patient treatment.

\emph{LSO intrinsic background.} The \iso{176}{Lu} content of LSO (2.6\% of
natural Lu) produces a \textbf{stationary $\sim$930--1000~cps random background}
independent of the signal. Because random coincidences scale as the square of the
singles rate, this intrinsic noise dominates at the low count rates of therapy
monitoring: the noise-to-signal ratio exceeded $0.2$ for true rates below
$\sim$4.8--5.7~kcps (about 1~kBq\,ml$^{-1}$, $\sim$1~MBq total). This is a
concrete penalty of an intrinsically radioactive scintillator in this
low-signal application --- the axis on which BGO and CsI, which carry no such
intrinsic activity, differ from LSO/LYSO. A separate finding, that PET range
verification is insensitive to CT artifacts from the titanium rods yet more
sensitive to the fluence ``shadowing'' behind them than a commercial pencil-beam
plan, is not pursued here.

\subsection*{2. Method and range metrics}

Range verification compares depth profiles of the measured and simulated
activity in the distal fall-off, at three definitions the authors label A, B, C:
\emph{method A} at the 20\% level of the last local maximum, \emph{method B} at
the 50\% level, and \emph{method C} (``shift method'') using the whole fall-off
region rather than a single point. The shift method is the more robust of the
three because it relies on the overall shape agreement rather than on a single
landmark, and it dominates the soft-tissue conclusions below.

\subsection*{3. Reproducibility (the enabling result)}

Two patients were imaged twice, one week apart, for the same treatment fraction.
The measured activity \emph{range} reproduced to $\sim$1~mm standard deviation
(mean absolute range deviation 1.14~mm and 0.125~mm for the two patients, shift
method), consistent with the 1~mm reproducibility seen earlier in phantoms.
Absolute activity values differed by up to $\pm30\%$ between the two scans, but
this is \emph{not} a timing effect: the MC, propagating the two different
irradiation durations and delays, predicts at most 3\% difference. The $\pm30\%$
therefore reflects statistical noise plus genuine day-to-day changes in tissue
composition and perfusion. \textbf{The 1~mm range reproducibility is what makes
the whole approach viable}; everything below is about the systematic offsets that
sit on top of it.

\subsection*{4. Biological washout}

Six head-and-neck patients (chosen for good immobilisation and rigid geometry, so
motion and co-registration are minimised) were analysed at positions sorted by
where the beam stops. The mean measured-vs-simulated range agreement is given in
Table~\ref{tab:washout}.

\begin{table}[h]
  \centering
  \caption{Mean agreement between measured and MC-simulated activity range,
    for the three stopping cases and the three range metrics
    (Knopf et al.~\cite{knopf2009}, Table~1). Case (i): beam stops in bone;
    case (ii): stops in soft tissue with the last activity maximum in bone;
    case (iii): stops in soft tissue with the last maximum also in soft tissue.}
  \label{tab:washout}
  \begin{tabular}{lccc}
    \toprule
    & \multicolumn{2}{c}{Pointwise (mm)} & Shift (mm) \\
    \cmidrule(lr){2-3}
    Case & 50\% (B) & 20\% (A) & (C) \\
    \midrule
    (i)   Bone                 & 2.5 & 1.2 & 2.4 \\
    (ii)  Bone / soft tissue   & 3.2 & 8.1 & 2.2 \\
    (iii) Soft tissue          & 6.8 & 3.9 & 4.3 \\
    \bottomrule
  \end{tabular}
\end{table}

\noindent
In bone (case i) all three metrics agree within $\sim$2.5~mm and verification
works, though individual positions can still deviate by up to 6.6~mm (50\%) /
3.6~mm (20\%). In soft tissue, pointwise verification breaks down because blood
perfusion degrades the measured distribution: the robust shift method reaches
4.3~mm in the worst case (iii) and 2.2~mm when the last maximum is anchored in
bone (case ii), but $27\%$ (case ii) / $33\%$ (case iii) of soft-tissue positions
are off by more than $\pm5$~mm. \textbf{Reliable mm-accurate verification in soft
tissue is therefore not feasible with the offline method.} The paper's often-quoted
``$\sim$4~mm washout discrepancy'' is precisely this case-(iii) soft-tissue number.
Two points matter for our reading: (a) this is a \emph{measured-vs-model bias}
driven by an imperfect washout model, not the intrinsic statistical $\sigma_R$;
and (b) the authors' proposed cure --- shorter delay via in-room PET to minimise
perfusion washout --- is the premise of the present scanner study.

\subsection*{5. Patient motion}

During the 30~min acquisition, breathing blurs the activity distribution. Lateral
conformity between measured and simulated distributions was up to 5~mm in
head-and-neck sites but degraded to as much as 30~mm in abdominopelvic sites,
dominant in the anterior--posterior direction. This excludes most thoracic and
abdominal tumours from mm-accurate verification with the 3D offline strategy,
independent of washout.

\subsection*{6. Uncertainty in the simulated activity (HU mapping)}

The MC activity depends on assigning elemental composition and washout parameters
to each voxel from its Hounsfield unit. Where tissues share an HU domain the
mapping is not correlative: bone marrow and soft tissue both fall in
$-20$ to $110$~HU, yet the Schneider conversion underestimates the marrow carbon
fraction by nearly 50\%, and since activation is dominated by the $(p,pn)$ channel
on \iso{12}{C} this underestimates marrow activity. Varying the HU conversion and
the washout parameters within plausible ranges moves the \emph{absolute} activity
by up to $\sim$35\% and the \emph{range} by $\sim$1.8--2.4~mm
(Table~\ref{tab:hu}), while the simulated \emph{dose} profiles are essentially
unaffected --- activity is far more sensitive to tissue parameters than dose is.

\begin{table}[h]
  \centering
  \caption{Activity-range differences induced by parameter changes in the MC
    (Knopf et al.~\cite{knopf2009}, Table~4): original vs adjusted HU conversion,
    and original vs adjusted washout parameters.}
  \label{tab:hu}
  \begin{tabular}{lccc}
    \toprule
    & 20\% (A) & 50\% (B) & Shift (C) \\
    \midrule
    $R_{\text{orig}}-R_{\text{adj.\ HU}}$ (mm)      & 1.80 & 2.41 & 1.81 \\
    $R_{\text{orig}}-R_{\text{adj.\ washout}}$ (mm) & 1.70 & 2.28 & 1.81 \\
    \bottomrule
  \end{tabular}
\end{table}

\noindent
The washout model itself is the space- and time-dependent weighting
$C_{\mathrm{bio}}(\mathbf{r},t) = M_s(\mathbf{r})\,
\exp[-\lambda_{\mathrm{bio}}(\mathbf{r})\,t]$ of the physical activity, with an
HU-based classification into low-, intermediate- and normally perfused tissue
(muscle: $M_s=0.55$, $T_{1/2,\mathrm{bio}}=3500$~s; brain: $M_s=0.35$,
$T_{1/2,\mathrm{bio}}=10\,000$~s), following Mizuno et al.~\cite{mizuno2003} and
neglecting any distinction between isotopes. These are the source constants used
in our own washout note.

\subsection*{7. Site-specific limitations}

Beyond the three general factors, offline verification is compromised for
particular sites by: \emph{beam arrangement} (opposed or small-angle fields blur
the distal edge of each field together --- feasible only for single or nearly
orthogonal beams); \emph{co-registration} (rigid CT fusion needs a rigid,
fully-imaged anatomy, which head/neck provides and eye/abdominopelvic do not);
and \emph{organ motion} (e.g.\ prostate patients voiding the bladder between
treatment and imaging changes the geometry irrecoverably).

\subsection*{8. Conclusion and relevance here}

In phantoms, where geometry and composition are known and motion and
co-registration are absent, the offline method reaches 1~mm range accuracy even
at 2~Gy-equivalent statistics (Sec.~1). In vivo, despite the same 1~mm range
reproducibility, offline PET/CT verifies the proton range to 1--2~mm only in
low-perfused, well-co-registered bony structures of head-and-neck patients.
Washout limits \emph{where} within a field verification is possible ($\sim$4~mm
in soft tissue); motion limits \emph{which whole sites} qualify (up to 3~cm in
the abdomen); and the MC activity is itself limited to $\sim$35\% in absolute
value and $\sim$2~mm in range by HU-to-tissue mapping. The authors argue all of
these are ``conquerable'', principally by \textbf{in-room PET} --- reducing the
treatment-to-scan gap to $\sim$1~min (their movable NeuroPET) suppresses washout
and motion, and removes the repositioning that drives co-registration error.
That prescription --- short delay, in-room acquisition --- is exactly the regime
the present $\sigma_R$-vs-dose study operates in; the LSO intrinsic-background
result (Sec.~1) is a quoted argument on the scintillator-choice axis that
separates BGO and CsI from LSO/LYSO; and the washout parameterisation summarised
in Sec.~6 is the one our downstream washout model adopts.

\begin{thebibliography}{9}

\bibitem{parodi2007a}
K.~Parodi, H.~Paganetti, E.~Cascio, J.~B.~Flanz, A.~A.~Bonab, N.~M.~Alpert,
K.~Lohmann and T.~Bortfeld,
``PET/CT imaging for treatment verification after proton therapy: a study with
plastic phantoms and metallic implants,''
\textit{Med.\ Phys.} \textbf{34} (2007) 419--435.

\bibitem{knopf2009}
A.~Knopf, K.~Parodi, T.~Bortfeld, H.~A.~Shih and H.~Paganetti,
``Systematic analysis of biological and physical limitations of proton beam
range verification with offline PET/CT scans,''
\textit{Phys.\ Med.\ Biol.} \textbf{54} (2009) 4477--4495.

\bibitem{mizuno2003}
H.~Mizuno, T.~Tomitani, M.~Kanazawa et al.,
``Washout measurements of radioisotope implanted by radioactive beams in the
rabbit,''
\textit{Phys.\ Med.\ Biol.} \textbf{48} (2003) 2269--2281.

\end{thebibliography}

\end{document}
